\subsection{A second executed table, with fundamental group $C_{23}$}
\label{D:C23-final-table-section}

\begin{theorem}\label{D:C23-final-table}
There is a connected finite category $\mathcal W_{23}$ with
\[
\begin{split}
 |\operatorname{Ob}\mathcal W_{23}|
   &=1\,021\,805\,843\,373\,670\,056,\\
 |\operatorname{Mor}\mathcal W_{23}|
   &=703\,363\,835\,733\,950\,740\,370\,112.
\end{split}
\]
Its fundamental group is $C_{23}$. Its classifying space is finitely
dominated, admits a CW model of dimension at most eight, and has
nonzero Wall finiteness obstruction. In particular the $C_{23}$ case
of D2 is negative, with a fixed finite composition table.
The constructor is \path{code/d_c23_wall_category.py}:
\texttt{build\_c23\_category()['Final']} returns the table with
the operations \texttt{source}, \texttt{target}, \texttt{identity},
and \texttt{compose}. Its object and arrow sets are the consecutive
integer intervals with the displayed cardinalities.
\end{theorem}

\begin{proof}
We specify the new finite inputs and the arithmetic indexing, then
apply the fully typed construction and signed chain calculation in
Theorem~\ref{D:Q8-final-table}. This instance was executed separately;
its definition does not alter the quaternion table.

\paragraph{The shorter integral lift.}
Set $R=\Z C_{23}$ and use the elements $N,\omega$ of
Proposition~\ref{D:C23-projector}. Define
\[
 e=\begin{pmatrix}
 4&\omega-1-N\\
 2\omega-N&-3
 \end{pmatrix},\qquad q=1-e.
\]
The three exact group-ring identities
\[
 N^2=23N,\qquad N\omega=12N,\qquad
 \omega^2-\omega+6=6N
\]
give
\[
\begin{split}
 (\omega-1-N)(2\omega-N)
 &=2\omega^2-2\omega-3N\omega+N+N^2\\
 &=-12.
\end{split}
\]
Thus $\operatorname{tr}e=1$, $\det e=0$, and the two-by-two
Cayley--Hamilton identity over the commutative ring $R$ gives
$e^2=e$. It follows that $q^2=q$ and $eq=qe=0$.
The augmentation of $e$ is
$\left(\begin{smallmatrix}4&-12\\1&-3\end{smallmatrix}\right)$.
Each of $e,q$ has exactly 42 signed group-labelled matrix units
when coefficients are expanded with their multiplicities. The full
coefficient arrays and the four projector identities are checked
independently by \path{code/d_c23_short_lift.py} and by the final
table constructor.

We verify that changing the augmentation lift preserves the required
nonvanishing conclusion. Under $R\to\mathcal O_L=\Z[\zeta_{23}]$,
the norm element $N$ vanishes. The image of $e$ is therefore exactly
\[
 \bar e=\begin{pmatrix}4&\omega-1\\2\omega&-3\end{pmatrix},
\]
the same cyclotomic projector used in
Proposition~\ref{D:C23-projector}. Let
$K=\Q(\omega)=\Q(\sqrt{-23})$ and
$I=(2,\omega)\subset\mathcal O_K=\Z[\omega]$.
The ideal has norm two and is not principal: the norm form
\[
 a^2+ab+6b^2=(a+b/2)^2+23b^2/4
\]
does not take the value two on integers. The element $\omega-2$
has norm eight and is absent from the conjugate prime
$(2,\omega-1)$, so $I^3=(\omega-2)$. Thus $[I]$ has exact order
three. Since $[L:K]=11$, taking ideal norm shows that
$I\mathcal O_L$ cannot be principal: otherwise $I^{11}$ would
be principal, contrary to its order.
Writing
\[
 p=\binom{2}{\omega},\qquad
 t=\bigl(2,(\omega-1)/2\bigr)
\]
gives $tp=1$ and $pt=\bar e$. Consequently
$\operatorname{im}\bar e\cong(I\mathcal O_L)^{-1}$.
Its determinant line is nontrivial, so it is not stably free.
Extension of scalars would send a zero reduced projective class
to zero. Therefore
\[
 [\operatorname{im}e]\ne0\quad\text{in }\widetilde K_0(R).
\]
This argument asserts nonvanishing for the new lift; equality with
the old integral projective class is not needed.

\paragraph{The finite base and all multiplication disks.}
Let $L$ be the join of two 23-gon circles. Number the first circle
by $0,\ldots,22$ and the second by $23,\ldots,45$, with their
cyclic edge orders. The group $C_{23}$ acts by simultaneous
addition modulo 23 in both layers. Since a nonidentity rotation
fixes no point in either circle, this action on
$|L|=S^1*S^1\cong S^3$ is free. For each nonempty face orbit
choose its lexicographically least face. Define $K_{23}$ by these
representatives and the arrows
\[
 (F,H,t):F\longrightarrow H\quad(t+F\subseteq H),
 \qquad
 (H,J,u)(F,H,t)=(F,J,u+t),
\]
with labels in $\Z/23$. There are 96 objects and 744 arrows;
its nondegenerate nerve counts are $(96,648,1104,552)$.
The face poset of $L$ is its free covering category. Hence
$BK_{23}\cong|\operatorname{sd}L|/C_{23}$ has fundamental group
$C_{23}$ and universal cover $S^3$.

For a label $a\in\{1,\ldots,22\}$ take the path
$(0,1,\ldots,a)$ in the first polygon. The identity label uses
the padded constant path $(0,0)$. For an inner label $h$ and
outer label $g$, concatenate the first path with its translate
$h+\operatorname{path}(g)$, then close it by the reverse of the
chosen path for $(h+g)\bmod23$. Cone this parameter polygon to
vertex 23 in the second layer. Every labelled triangle is a face
of the join, possibly with a collapsed boundary edge. This fixes
all 529 multiplication disks; their boundary lengths lie between
3 and 65. For an inclusion of parameter faces $F\subseteq H$,
choose the unique transport labels $t_F,t_H$ from their normalized
image faces and use the arrow label $t_F-t_H$. This supplies exact
functors to $K_{23}$ and exact boundary equalities. Their complete
native tables are saved in
\path{code/outputs/d_c23_coefficient_context_certificate.json.gz}.

Use $Y=NK_{23}\vee A_1^{\wedge5}\vee A_1^{\wedge5}$ and the same
signed elementary cospans, flat pinches, permutation grids, and
simultaneous cancellation cylinders as in
Theorem~\ref{D:Q8-final-table}. Expanding each of $eq$ and $qe$
gives 1764 terms. The deterministic pairing order is lexicographic
in the target index, group label, and source index, with occurrences
paired in their original order. For $eq$ this yields 405 opposite
pairs and 954 zero products; for $qe$ it yields 416 opposite pairs
and 932 zero products. Each specified nullhomotopy has 50 cospan
roofs and 49 actual comparison maps. All occurrences of a group
product use the corresponding fixed fan disk.

\paragraph{The arithmetic table and its realization.}
The finite simplicial normal forms and ordinal pullback rules are
unchanged. For $A_n$ use this fixed arithmetic order in degree $k$:
the basepoint, the $n-1$ constant interior vertices, the $n$ constant
edges, then the $2nk$ crossing strings ordered by edge, endpoint
side, and the number of initial vertices. Thus all ranks and
unranks are given by integer arithmetic. The parameter interval
$J_n$ lists first its vertices, then its edges; its arrows list
identities, left incidences, then right incidences. The optional
backend \path{code/d_arithmetic_intervals.py} implements precisely
these orders and the same composition functions. The displayed
constructor is run in a fresh Python interpreter, installing this
backend before constructing the new instance; its process-local
factory replacements are not mixed with a quaternion-table run.

Apply the finite simplex-category replacement $\mathcal O_8$ to
each coefficient space, then the ordinal, product, span, and
height-one constructors already specified in
Theorem~\ref{D:Q8-final-table}. Every coefficient roof has dimension
at most seven; the final relative nullhomotopy cylinders have
dimension at most eight. Thus the proved truncation equivalence
$B\mathcal O_8X\simeq|X|$ applies to every input. The same two
parity monoids, suspension comparison, tensor-zero contraction,
base-path collars, and intrinsic pinch produce the final finite
category. Its 3897 homogeneous type blocks have exactly
\[
\begin{split}
 \#\{\text{typed pairs}\}
  &=288\,986\,603\,572\,868\,508\,990\,808\,637\,052,\\
 \#\{\text{typed triples}\}
  &=109\,162\,136\,788\,728\,112\,368\,563\,487\,083\,799\,108.
\end{split}
\]
Every lookup is a finite recursion on these constructor types;
there is no subdivision search or unstated choice in the table.

\paragraph{Chains, domination, and obstruction.}
All maps are over $BK_{23}$, with their under-base identifications
retained as the explicit natural-transformation paths. Van Kampen
therefore gives $\pi_1(B\mathcal W_{23})=C_{23}$. The same parity
calculation gives relative chains $R^2$ in each degree at least
seven, and alternates $e,q$. After the first cone and its explicit
deformation retraction, the second attaching map has coordinates
\[
 x\longmapsto\bigl(\alpha_n(1-f_n)x,\,\beta_nf_nx\bigr),
 \qquad \alpha_n,\beta_n\in\{1,-1\}.
\]
Its inverse on $(1-f_n)R^2\oplus f_nR^2$ is
$(u,v)\mapsto\alpha_nu+\beta_nv$. Thus every positive relative
degree contracts, leaving $qR^2$ in degree seven. The coefficient
nullhomotopies have zero degree-five-to-six cellular homotopy
term, and the extra tensor-zero coherence factors through the
zero relative coefficient quotient, exactly as in the first
table. No additional matrix commutativity or sign assumption
enters the cone calculation.
Consequently
\[
 C_*(\widetilde{B\mathcal W_{23}})\simeq
 C_*(\widetilde{BK_{23}})\oplus qR^2[7],
\]
where the base complex has free ranks $(96,648,1104,552)$ in
degrees zero through three.

The finite seven-skeleton of the nerve dominates the whole space:
its inclusion is seven-connected, and all section-obstruction
groups in degrees at least eight vanish for every local coefficient
system by this bounded projective chain model. A mapping telescope
of the resulting homotopy idempotent on that finite skeleton gives
a CW model of dimension at most eight. Finally Wall's obstruction
is
\[
 w(B\mathcal W_{23})=-[qR^2]=[eR^2]\ne0.
\]
Wall's theorem therefore excludes a finite CW model.

\paragraph{Executed checks.}
All 529 native fan disks were checked, with 911\,881 composable
triples and 107\,272 exact boundary equations. The independent
elementary-coefficient checker covered all 4232 group/sign/index
cases and 338\,560 endpoint equations. Each complete nullhomotopy
passed 2940 selected endpoint/pullback equations. The final run
passed 558 base-retraction equations, 468 naturality equations,
and 15\,588 identity/associativity equations at the ends of all
3897 type blocks. The last checks sample the large table; the
universal constructor proofs establish its laws for every index.
The exact record is
\path{code/outputs/d_c23_wall_category_summary.json}.
\end{proof}
